% !TEX program = lualatex
% LuaLaTeX でコンパイル: latexmk -lualatex slash_commands_comparison.tex
\documentclass[a4paper,11pt]{ltjsarticle}

\usepackage[margin=25mm]{geometry}
\usepackage{booktabs}
\usepackage{tabularx}
\usepackage{array}
\usepackage{enumitem}
\usepackage{xcolor}
\usepackage[hidelinks]{hyperref}
\usepackage{xurl}

\newcolumntype{L}{>{\raggedright\arraybackslash}X}
\newcolumntype{P}[1]{>{\raggedright\arraybackslash}p{#1}}
\newcommand{\cmd}[1]{\texttt{#1}}
\newcommand{\home}{\textasciitilde}
\renewcommand{\arraystretch}{1.25}

\title{2026年における主要AI開発エージェントの\\アーキテクチャとSlash Commandsの徹底解剖：\\Antigravity、Claude Code、Codexの比較分析}
\author{}
\date{2026年9月}

\begin{document}
\maketitle

\begin{abstract}
ソフトウェア開発におけるAIの役割は、エディタ内の単一ファイル補完から、リポジトリ全体を探索し計画・実行・検証を自律的に行うエージェントへと移行した。
本稿では、こうしたエージェントとの対話において決定論的かつ再現性の高い操作を提供する「スラッシュコマンド」に着目し、GoogleのAntigravity、AnthropicのClaude Code、OpenAIのCodex CLIの3ツールについて、組み込みコマンドの体系、カスタムコマンドの構築メカニズム、基盤アーキテクチャの差異を比較分析する。
さらに、スラッシュコマンドを支える技術基盤としてModel Context Protocol（MCP）とコンテキストエンジニアリングを整理し、エージェントの自律性がもたらす間接的プロンプトインジェクションや「GuardFall」と呼ばれるセキュリティ上の脅威と、その対策について論じる。
\end{abstract}

\noindent\textbf{キーワード：} AIコーディングエージェント，スラッシュコマンド，Model Context Protocol，コンテキストエンジニアリング，プロンプトインジェクション

\section{導入：AIソフトウェア開発エージェントにおける操作パラダイムの移行}

ソフトウェアエンジニアリングの領域において、人工知能の役割は劇的な変化を遂げている。数年前まで主流であった「エディタ内での単一ファイルのコード補完（Copilot型）」から、現在ではAIが自律的にリポジトリ全体を探索し、計画を立案し、ターミナルコマンドを実行し、テストを修正する「自律型エージェント」へと移行している\cite{r1}。この高度な自律性を持つエージェントとの対話において、自然言語の曖昧さを排除し、開発者が決定論的かつ再現性の高い操作を行うための標準的なインターフェースとして定着したのが「Slash Commands（スラッシュコマンド）」である\cite{r3}。

スラッシュコマンドは、チャットインターフェースの入力プロンプトにおいて \cmd{/}（スラッシュ）から始まる特定の文字列を入力することで、AIに対して事前に定義された複雑なワークフローの実行、認知モードの切り替え、または特定のコンテキスト（文脈）のロードを強制するメカニズムである\cite{r6}。これにより、開発者は「コンテキストの圧縮」「別スレッドでの並行検証」「特定のセキュリティレビュー」といった高度な操作を、自然言語で長々と説明することなく、瞬時に呼び出すことが可能となる\cite{r8}。

本稿では、2026年現在、世界の開発現場を牽引している3つの主要な自律型AIコーディングエージェント——Googleの「Antigravity（旧Gemini CLI）」、Anthropicの「Claude Code」、およびOpenAIの「Codex CLI」——に焦点を当て、各ツールに実装されているスラッシュコマンドの内部仕様、カスタムコマンドの構築メカニズム、基盤となるアーキテクチャの差異、そして自律型エージェント特有のセキュリティ脅威について網羅的かつ深層的な分析を行う。

\section{スラッシュコマンドを支える次世代プロトコルとアーキテクチャ}

各ツールの具体的なコマンド仕様を分析する前に、最新のスラッシュコマンドを支える技術的基盤と、それが解決しようとしている大規模言語モデル（LLM）特有の課題について理解することが不可欠である。

\subsection{Model Context Protocol (MCP) によるプロンプトの標準化と露出}

2024年後半にAnthropicによって提唱され、2026年現在ではAI業界全体のオープンスタンダードとして広く採用されているModel Context Protocol（MCP）は、AIエージェント（クライアント）が外部のリソース、ツール、およびプロンプトにアクセスするためのJSON-RPCベースの標準規格である\cite{r3}。MCPはしばしば「AIのためのUSB-C」と形容され、エージェントと外部データソース間の通信を完全に標準化する\cite{r11}。

MCP仕様において、スラッシュコマンドは主にMCPの「Prompts（プロンプト）」プリミティブのユーザーインターフェースとしての側面を持つ\cite{r4}。MCPサーバーは、動的な引数を受け取り、リソースからのコンテキストを包含し、特定のワークフローを誘導する再利用可能なプロンプトテンプレートを定義できる\cite{r4}。開発者がチャットウィンドウで \cmd{/} を入力した際に表示されるコマンド群の一部は、ハードコードされたものではなく、実際には \cmd{prompts/list} エンドポイントを通じて接続中のMCPサーバーから動的に取得・露出されたものである\cite{r3}。これにより、例えばSlackやBigQuery、あるいは社内専用のデータ基盤と連携するMCPサーバーを接続した瞬間、それらを操作するための専用スラッシュコマンドが自動的にターミナル上で利用可能になるという、極めて拡張性の高いエコシステムが形成されている\cite{r3}。

\subsection{トークン爆発（Token Explosion）の抑制とコンテキストエンジニアリング}

スラッシュコマンドが現代のAIエージェントにおいて重用される第二にして最大の理由は、LLMの推論効率の最適化とコンテキストウィンドウの枯渇防止にある。多数のツール（関数スキーマ）をシステムプロンプトに常駐させると、1つのツール定義につき150〜400トークンを消費し、20個のツールを読み込むだけで数千から数万トークンが消費される「トークン爆発（Token Explosion）」が発生する\cite{r15}。この状態に陥ると、モデルの注意機構（Attention Mechanism）が散漫になり、ツールの誤選択、ハルシネーション、あるいはコンテキストの忘却（Context Rot）が引き起こされる\cite{r15}。

スラッシュコマンドは、この問題に対する「遅延読み込み（Lazy Loading）」と「スコープ制御」の実践的な解決策を提供する\cite{r15}。ユーザーが \cmd{/review} や \cmd{/plan} といった特定のスラッシュコマンドを呼び出した瞬間にのみ、そのタスクに必要な専用のツール群（例：コードの差分読み込みツールやGit操作ツール）と指示文がコンテキストに注入される。これにより、エージェントは常にそのターンの目的に対して最小十分なコンテキスト（Minimal Sufficiency）のみを保持してタスクを実行でき、推論の精度向上とレイテンシの劇的な削減が実現される\cite{r15}。

\section{Claude Code：深層推論と柔軟なMarkdownベースの拡張性}

Anthropicが提供するClaude Codeは、深い推論（Deep Reasoning）に最適化されたCLIベースのエージェントツールであり、Claude ProやMaxなどのサブスクリプションプランに内包される形で提供されている\cite{r18}。最大100万トークンという巨大なコンテキストウィンドウと、OSレベルでのサンドボックス強制、破壊的コマンドに対する厳格な承認フローを備えており、大規模なリファクタリングや複雑なバグ調査において高い能力を発揮する\cite{r18}。

\subsection{組み込みコマンドの体系とセッション管理}

Claude Codeには、2026年時点で60種類以上の組み込みスラッシュコマンドが搭載されており、日常的な開発操作のほぼすべてがコマンド化されている\cite{r9}。これらのコマンドは、機能の抽象度や目的によって高度に分類されている（表\ref{tab:claude}）。

\begin{table}[htbp]
\centering
\caption{Claude Codeの主要な組み込みスラッシュコマンド}
\label{tab:claude}
\small
\begin{tabularx}{\textwidth}{P{30mm}P{24mm}L}
\toprule
コマンド & カテゴリ & 役割とメカニズム \\
\midrule
\cmd{/compact [指示]} & コンテキスト管理 & 長期化したセッションの会話履歴を要約して圧縮する。トークン消費を抑えつつ、直前の文脈や重要な決定事項を維持するために不可欠なコマンドである\cite{r9}。 \\
\cmd{/branch [名前]} & コンテキスト管理 & 現在のセッションの状態から別のバックグラウンドスレッド（ブランチ）をフォークし、メインのターミナルを解放する。長時間かかるテストや並行検証に用いられる（エイリアス：\cmd{/fork}, \cmd{/bg}）\cite{r9}。 \\
\cmd{/clear} & コンテキスト管理 & 会話履歴を完全にクリアし、コンテキストウィンドウを初期状態にリセットする。前の状態には \cmd{/resume} を用いて復帰可能である\cite{r6}。 \\
\cmd{/rewind} & コンテキスト管理 & エージェントが誤ったコード修正を行った際などに、会話やコードの状態を以前のチェックポイント（時点）に巻き戻す\cite{r9}。 \\
\cmd{/review [PR]} & コード・レビュー & 指定されたプルリクエスト（PR）をローカルセッションにロードし、マルチエージェントによる詳細なコードレビューを実行する。より深いクラウドサンドボックス上でのレビューには \cmd{/ultrareview} が用いられる\cite{r9}。 \\
\cmd{/security-review} & コード・レビュー & 現在のブランチの未コミット変更を、セキュリティの観点から深層分析する\cite{r9}。 \\
\cmd{/plan [説明]} & バンドルスキル & 即座に実装に入るのではなく、コードベース全体を調査し、要件を定義し、実装計画を立案する「プランモード」に移行する。説明を引数として渡すことで、特定のタスクの計画から開始できる\cite{r9}。 \\
\cmd{/batch <指示>} & バンドルスキル & コードベース全体に対する大規模な変更を、5〜30個の独立したユニットに分解し、複数の並列処理エージェントに委譲して一括処理を行う\cite{r9}。 \\
\cmd{/simplify} & バンドルスキル & 変更されたファイルを3つのエージェントによって並列でレビューさせ、コードの品質や効率に関する問題を特定し、単純化・修正する\cite{r9}。 \\
\cmd{/btw <質問>} & ユーティリティ & メインの会話コンテキストを汚染することなく、バックグラウンドでサッと簡単な質問（By the way）を実行する\cite{r8}。 \\
\cmd{/hooks} & 設定と管理 & セッション終了時やツール使用前後に自動実行される「Hooks」の設定を確認・管理する\cite{r9}。 \\
\bottomrule
\end{tabularx}
\end{table}

\subsection{カスタムコマンドの構築メカニズム}

Claude Codeのアーキテクチャにおける最大の強みは、開発者がプロジェクト固有のカスタムスラッシュコマンドを極めて直感的かつ容易に構築できる拡張性にある。開発者は、プロジェクトのルートディレクトリにある \cmd{.claude/commands/}、またはユーザーグローバルの \cmd{\home/.claude/commands/} 配下にMarkdownファイル（\cmd{.md}）を配置するだけでよい。このファイル名から拡張子を除いた部分が、そのまま独自のスラッシュコマンド名として認識される（例：\cmd{fix-issue.md} は \cmd{/fix-issue} となる）\cite{r6}。

カスタムコマンドのMarkdownファイルは、単なるプロンプトのテンプレートを超えた、強力なメタデータ制御とコンテキスト動的注入機能を備えている。

\begin{enumerate}[label=(\arabic*)]
\item \textbf{YAMLフロントマターによるメタデータ定義}：ファイルの先頭に記述されるYAMLフロントマター（Frontmatter）を用いて、メタデータを定義する。ここで \cmd{allowed-tools} 属性を指定することで、そのコマンドの実行中にAIが使用できるツールを厳密に制限できる（例：\cmd{Bash(git:*)}, \cmd{Bash(npm:*)}, \cmd{Read(*.md)}）\cite{r7}。これにより、特定の調査コマンド実行時にエージェントが暴走して無関係なファイルを書き換えるといったリスクを構造的に排除できる\cite{r7}。また、\cmd{argument-hint} を定義することで、CLI上で引数の入力補助を表示させることが可能である\cite{r9}。
\item \textbf{引数の動的展開}：コマンド呼び出し時にユーザーが入力した文字列は、Markdownファイル内の \cmd{\$ARGUMENTS} という変数に動的に展開される\cite{r6}。
\item \textbf{シェル実行の事前処理}：ファイル内の行頭に \cmd{!}（感嘆符）を記述すると、プロンプトがLLMに送信される前に指定されたBashコマンドがローカル環境で実行され、その標準出力結果がスラッシュコマンドのコンテキストとして動的に挿入される\cite{r7}。例えば、\cmd{!~npx textlint \$ARGUMENTS.md} と記述すれば、AIに指示を出す前に静的解析ツールを実行し、そのエラー結果をコンテキストに含めた上で修正を依頼するといった自律的なワークフローが1つのコマンドで完結する\cite{r7}。
\item \textbf{ファイル参照}：行頭に \cmd{@} を付けてファイルパスを指定すると、対象ファイルの内容がそのままプロンプトのコンテキストとして展開され、参照コードとして活用される\cite{r7}。
\end{enumerate}

2026年の最新の運用動向として、ユーザーが任意のタイミングで手動で呼び出す旧来の \cmd{.claude/commands/*.md} 形式に加え、AIが状況（コンテキスト）に応じて自律的に使用タイミングを判断し自動起動する機能を持つ \cmd{.claude/skills/SKILL.md} への移行が公式に推奨されている。これにより、開発体験は「呼ばれたら動く」コマンド駆動から、「勝手に動く」スキル・フック駆動へとパラダイムシフトが進行している\cite{r9}。

\section{Google Antigravity：マルチエージェントオーケストレーションと非同期実行}

Googleが提供する「Antigravity」は、2026年6月18日をもってサービス終了したGemini CLIの完全なる後継プラットフォームとして、Google I/O 2026で発表されたAIネイティブ開発環境である\cite{r23}。Antigravityは、Go言語で完全に再設計され、極めて高い処理速度と軽量性を誇るターミナル用の「Antigravity CLI」と、統合的なビジュアルエディタである「Antigravity 2.0 (GUI)」という2つのサーフェスから構成される\cite{r23}。両者は全く同一のコアエージェントハーネスを共有しており、設定、セキュリティ権限、さらには進行中の会話コンテキストに至るまで、インターフェース間でシームレスに同期・移行が可能である\cite{r26}。

\subsection{組み込みコマンドの体系と推論モードの制御}

Antigravityのスラッシュコマンドは、Claude Codeのそれと比較して、バックグラウンドでの並列処理と、タスクの複雑さに応じた「推論モード」の動的な切り替えに強く特化している（表\ref{tab:antigravity}）。

\begin{table}[htbp]
\centering
\caption{Antigravityの主要な組み込みスラッシュコマンド}
\label{tab:antigravity}
\small
\begin{tabularx}{\textwidth}{P{30mm}P{18mm}LP{18mm}}
\toprule
コマンド & カテゴリ & 役割とメカニズム & 対象タイム\newline ホライズン \\
\midrule
\cmd{/boost} & 推論と自律性 & 複雑なバグ、競合状態、アルゴリズムの最適化といった高度な課題に対して、「オーケストレーター」「ディープコーダー」「検証ワーカー」からなる3層のマルチエージェント推論階層をオンデマンドで起動する。これらは独立して仮説を立て、テストスイートを実行し、結果を独自に検証してからメインスレッドに返却する\cite{r8}。 & 秒〜数時間 \\
\cmd{/teamwork-preview} & 推論と自律性 & リポジトリ規模のマイグレーションや長期研究プロジェクトのために、自律的なエージェントチームを編成する。事前のスコープ定義インタビューを行い、マイルストーンに沿ってロードマップを作成し、並列作業を実行する（有料プラン限定機能）\cite{r8}。 & 数時間〜数日 \\
\cmd{/goal} & 推論と自律性 & ユーザーのターンごとの確認（Turn-by-turn confirmation pauses）をスキップし、指定された目標（例：特定のパッケージ内の全ユニットテストをパスさせる）が達成されるまで、連続的かつ自律的に実行を続ける\cite{r8}。 & 分〜数時間 \\
\cmd{/plan} & 計画 & 直接コードを編集する前に、コードベースを調査し、要件定義のインタビューを行い、レビュー可能な計画書（アーティファクト）を作成する。複雑なタスクのスコープを事前に整理するために用いられる\cite{r8}。 & 分 \\
\cmd{/grill-me} & 計画 & 曖昧な要件に対して、エッジケースや制約条件を明確にするためのステップバイステップの対話的インタビューをAI側からユーザーに対して実施する\cite{r8}。 & 分 \\
\cmd{/learn} & カスタマイズ & 直近のセッションでの修正履歴やユーザーからのフィードバックパターンを分析し、将来のセッションで永続的に適用されるプロジェクトの「ルール」や「スキル」として蒸留・保存する\cite{r8}。 & 即時 \\
\cmd{/schedule} & 自動化 & 指定されたタスクを、一回限りのタイマー、またはcron式のスケジュール（例：毎朝特定のブランチをクリーンアップする）に従ってバックグラウンドで自律実行させる\cite{r8}。 & スケジュール指定 \\
\cmd{/browser} & ツール & サンドボックス化されたChromeブラウザのサブエージェントを起動し、Web上の最新リサーチや、ローカルで立ち上げたUIのライブレンダリング・レイアウト検証を実行させる\cite{r8}。 & 分 \\
\cmd{/title} & ウィンドウ制御 & ターミナルウィンドウのタイトルバー機能を切り替え、アクティブなモデル、ワークスペース、エージェントのステータスを動的に表示するよう構成する\cite{r29}。 & 即時 \\
\cmd{/voice} & インプット & 音声ディクテーションを有効化し、タイピングではなく音声によってエージェントへの指示を入力する\cite{r30}。 & 即時 \\
\bottomrule
\end{tabularx}
\end{table}

Antigravity CLIの最も革新的なUX（ユーザーエクスペリエンス）の一つが、ターミナル上で非同期に実行されるサブエージェント群を監視・管理する機能である。プロンプト内で \cmd{/agents} と入力すると、専用のサブエージェント管理パネルが開き、実行中、完了、または強制終了されたサブエージェントのステータスと、現在実行中のステップが一覧表示される\cite{r28}。開発者はメインの会話フローを中断することなく、\cmd{ctrl+j} を押下して承認待ちのサブエージェントの詳細ビューにジャンプしたり、「Fast Path Alerts」機能を利用して \cmd{ctrl+k} でツールの実行許可を即座に承認したりすることが可能である\cite{r28}。

\subsection{カスタムコマンドと外部統合（TOMLとMCP）}

Antigravityにおけるカスタムスラッシュコマンドの定義方法は、Claude CodeのMarkdown形式とは異なり、TOML形式（\cmd{.toml}）のファイルとして構成される。プロジェクトスコープの場合は \cmd{.gemini/antigravity-cli/} 配下（旧 \cmd{.gemini/commands/}）、ユーザースコープの場合は \cmd{\home/.gemini/commands/} に配置される\cite{r18}。

このTOMLファイル内では、コマンドの \cmd{description}（説明）と \cmd{prompt}（指示内容）を定義し、ユーザーからの入力は \cmd{\{\{args\}\}} 変数としてプロンプト内に展開される\cite{r32}。さらに特筆すべきは、Antigravity CLIがFastMCPなどのMCPサーバーとネイティブに深く統合されており、設定されたMCPサーバーが提供するプロンプトを自動的に解析し、スラッシュコマンドとしてターミナルに露出させる機能を持っている点である\cite{r26}。これにより、例えばDataRobotなどのサードパーティプラットフォームが提供する専用のスキルセットプラグインをインストールするだけで、\cmd{/datarobot-skills:datarobot-agent-assist} といった高度な専用スラッシュコマンドが、自作することなく即座に利用可能となる\cite{r24}。これらの設定やキーバインドは、すべて \cmd{\home/.gemini/antigravity-cli/settings.json} および \cmd{keybindings.json} で詳細にカスタマイズ可能である\cite{r28}。

\section{OpenAI Codex (Codex CLI)：厳密な階層型コンテキストと非対話型自動化}

OpenAIが提供する「Codex CLI」は、ChatGPTのサブスクリプション（Plus, Pro, Business, Enterprise等）にバンドルされる形で提供される、ターミナル常駐型のAIコーディングエージェントである\cite{r18}。Claude CodeやAntigravityが自然言語による「会話型のインターフェース」に重きを置いているのに対し、Codex CLIはより「ソフトウェア・エンジニアリング・エージェント」としての性格が強く、厳密なサンドボックス制御、非対話型の自動化（\cmd{codex exec} によるCI/CDパイプラインへの組み込み）、そしてプロジェクト規約の徹底に特徴がある\cite{r18}。

\subsection{組み込みコマンドとセッション制御}

Codex CLIのスラッシュコマンドは、対話の管理、状態の監視、およびワークフローの厳格な切り替えに重点を置いている（表\ref{tab:codex}）。

\begin{table}[htbp]
\centering
\caption{Codex CLIの主要な組み込みスラッシュコマンド}
\label{tab:codex}
\small
\begin{tabularx}{\textwidth}{P{38mm}L}
\toprule
コマンド & 役割とメカニズム \\
\midrule
\cmd{/permissions} & セッション中にコマンド実行前の承認挙動やサンドボックスの実行権限を動的に変更する（旧 \cmd{/approvals}）\cite{r37}。 \\
\cmd{/status} & 現在のアクティブなモデル、承認モード、トークン使用量、コスト見積もり、Weeklyリミットの消化状況、およびセッションの詳細などのシステムステータスを一覧表示する\cite{r36}。 \\
\cmd{/usage} & トークンのアクティビティ、レート制限、および利用可能な使用量リセットクレジットの状況を表示する\cite{r41}。 \\
\cmd{/compact} & 長期化し蓄積された会話履歴を要約し、コンテキストウィンドウの消費を節約する\cite{r36}。 \\
\cmd{/side} & 現在のメインスレッドのコンテキストを維持したまま、一時的なサイドチャット（別スレッド）を開始する。本筋のコードを汚さずにアーキテクチャの相談や特定の関数の動作検証を行う際に極めて有効である\cite{r37}。 \\
\cmd{/review} & 対話モードを離れ、現在の作業ツリーの差分（変更点）に対する非対話型のコードレビューを実行し、バグやリグレッションを指摘する\cite{r37}。 \\
\cmd{/init} & カレントディレクトリを自動探索し、プロジェクトの概要や技術スタックを記述した \cmd{AGENTS.md} の足場（スキャフォールド）を自動生成する\cite{r36}。 \\
\cmd{/model} & セッションの途中でアクティブなモデルと、推論の深さ（Reasoning effort: low, medium, high）を動的に切り替える\cite{r36}。 \\
\cmd{/sandbox-add-read-dir} & サンドボックス環境に対して、指定した別のディレクトリへの読み取りアクセス権を動的に付与する\cite{r39}。 \\
\cmd{/export} & 現在のターミナル（TUI）上の会話全体をMarkdown形式でエクスポートし、ファイルに保存するかクリップボードにコピーする\cite{r37}。 \\
\cmd{/cd} / \cmd{/pwd} & エージェントの作業ディレクトリをセッション内で動的に変更・表示する（v0.149.0以降の機能）\cite{r37}。 \\
\bottomrule
\end{tabularx}
\end{table}

\subsection{アーキテクチャの中核：階層型 AGENTS.md パラダイム}

Codex CLIのアーキテクチャを語る上で欠かせないのが、コンテキスト管理の絶対的な基盤となる \cmd{AGENTS.md} という標準仕様である。これはAgentic AI Foundationの管理下で60,000以上のプロジェクトで採用されているオープンフォーマットであり、AIエージェントに対してプロジェクト固有のコーディング規約、テスト環境の仕様、アーキテクチャの制約、そして絶対的な禁止事項を永続的に定義する「エージェント向けのREADME」として機能する\cite{r43}。

Codexは、新たなタスクを実行する際、必ず以下の順序で階層的に \cmd{AGENTS.md} を探索し、それらを結合（マージ）してプロンプトの最上流コンテキストに注入する\cite{r46}。

\begin{enumerate}
\item \textbf{グローバルスコープ}：ユーザーのホームディレクトリ（\cmd{\home/.codex/AGENTS.md}）。個人のコーディングの好みや、全プロジェクトに共通して適用すべき普遍的なルール（例：常に特定のパッケージマネージャーを優先する等）を定義する\cite{r45}。
\item \textbf{プロジェクトルート}：リポジトリ直下の \cmd{AGENTS.md}。プロジェクト全体のアーキテクチャ設計、CI/CD規則、フレームワークのバージョン制約など、チーム全体で共有すべき設定を記述する\cite{r45}。
\item \textbf{サブディレクトリ}：現在の作業ディレクトリに至るまでの各階層の \cmd{AGENTS.md} または \cmd{AGENTS.override.md}。特定のサービスやモジュール（例：\cmd{frontend/} や \cmd{services/payments/}）固有のローカルなルールを定義する\cite{r45}。
\end{enumerate}

重要なアーキテクチャ上の特性として、Codexはこれらのファイルをルートからカレントディレクトリに向かって順番に読み込み、結合する。プロンプト内では後に読み込まれたファイル（より深い階層のファイル）の指示が優先されるため、モジュール固有のローカルルールがプロジェクト全体のグローバルルールを安全に上書き（オーバーライド）する設計となっている\cite{r46}。

AIの挙動が不安定になること（ハルシネーションや規約違反）を防ぐため、\cmd{AGENTS.md} には自然言語の曖昧な散文ではなく、「検証可能なコマンド（例：\cmd{npm run lint} を実行し、終了コードが0であること）」や「Doneの定義（Definition of Done）」を厳格に記述することがベストプラクティスとされている\cite{r43}。設定ファイルである \cmd{\home/.codex/config.toml} では、使用モデルの指定に加え、ファイル書き込みの制限レベル（\cmd{workspace-write} 等のサンドボックスモード）や、プロファイルごとの細かな挙動制御が可能である\cite{r36}。また、CodexはMCPクライアントとしてだけでなく、自身がMCPサーバーとして機能する能力も備えており、他のツールから呼び出されるバックエンドエージェントとしても動作可能である\cite{r18}。

\section{3大プラットフォームの多角的な比較分析}

Antigravity、Claude Code、Codexは表面的には類似した「AIコーディングエージェント」という機能を提供するが、そのターゲット層、価格モデル、コンテキストウィンドウの設計、およびサブエージェントの処理方式において、明確な設計思想の差異が存在する\cite{r1}（表\ref{tab:compare}）。

\begin{table}[htbp]
\centering
\caption{3大プラットフォームの比較}
\label{tab:compare}
\footnotesize
\begin{tabularx}{\textwidth}{P{22mm}LLL}
\toprule
比較項目 & Google Antigravity & Claude Code & OpenAI Codex CLI \\
\midrule
価格モデルと利用制限 & 基本無料枠あり（制限超過後はフラッシュモデル等の軽量モデルへフォールバック）。GCP連携重視\cite{r18}。無料枠のクォータ上限に達しやすい\cite{r51}。 & Claude Pro/Maxサブスクリプション（月額\$20〜）に内包。API利用枠を共有するため、使用量制限は比較的厳しい\cite{r19}。 & ChatGPTサブスク（Plus/Pro/Biz等）に内包。高リミットで許容範囲が広く、長時間の作業に適する\cite{r18}。 \\
コンテキストウィンドウ & Gemini 1.5 Pro等の統合により最大約100万トークンの広大なウィンドウ\cite{r18}。 & 20万〜100万トークン。深層推論と長期的なコンテキスト保持に極めて優れる\cite{r18}。 & CLI本体では約40万トークン（API経由で最大100万まで拡張可）\cite{r18}。 \\
サブエージェント方式 & ターミナル上の専用パネル（\cmd{/agents}）による非同期監視。メインスレッドをブロックせず並行作業を委譲\cite{r18}。 & ネイティブな深い推論による自律的なタスク分解とサブワーカーの呼び出し\cite{r18}。 & 実験的なMulti-agents機能。環境変数で有効化し、ファンアウトと結果の集約によるサブタスク実行\cite{r18}。 \\
セキュリティ・サンドボックス & 軽量なOSネイティブ分離（Linuxのnsjail、macOSのsandbox-exec等）をゼロスタートアップ時間で提供\cite{r28}。 & OS強制サンドボックスに加え、破壊的コマンドに対する厳格な組み込み承認フロー（Bash/Zsh対応）\cite{r18}。 & \cmd{workspace-write} などの明示的な名前付きポリシー。隔離環境での実行（\cmd{codex exec}）に強み\cite{r18}。 \\
カスタムコマンド・拡張形式 & \cmd{.toml} 形式（プロンプトと設定を定義）およびMCPサーバーからの自動露出\cite{r32}。 & \cmd{.md} 形式（YAMLフロントマターによるツール制限、Bashコマンド実行を含む）\cite{r7}。 & \cmd{AGENTS.md} を軸とした階層的コンテキスト定義と設定ファイル（\cmd{config.toml}）\cite{r44}。 \\
\bottomrule
\end{tabularx}
\end{table}

\subsection{実践的ユースケースに基づく評価と開発者の選定基準}

実業務における評価データや開発者コミュニティのフィードバックに基づくと、これら3つのツールの選択は、開発チームのスタイルやタスクの性質に強く依存する\cite{r19}。

Claude Codeは、事前のプロンプト指示が少なくてもリポジトリ全体のコンテキストを把握し、文脈から最適解を導き出す「探求的・推論的」なタスクに圧倒的に優れている\cite{r19}。しかし、些細な修正依頼に対して広範なコード変更を提案する「オーバーエンジニアリング」の傾向があり、さらにトークンリミットの消費が早いため、週の半ばで制限に到達してしまうケースが報告されている\cite{r19}。

一方、Codex CLIは、仕様が明確に定義されたタスク（特に \cmd{AGENTS.md} によって制約が厳格にコントロールされた環境下）において、開発者の指示に忠実に従う実行精度の高さが評価されている\cite{r19}。また、サブスクリプションの利用枠（クォータ）が比較的大きく、長時間の非同期タスクやCI/CDパイプラインへのヘッドレスな組み込みに最適である\cite{r19}。

Antigravityはこれらの中間に位置する。GUI（Antigravity 2.0）とCLIのシームレスな移行、およびGoogle Cloudのエコシステム（例えばSpannerマイグレーションタスクの自動化など）との統合において無類の強みを発揮する\cite{r23}。初期導入のハードルが低く初学者にも推奨されるが、無料枠のクォータ上限に達しやすく、ヘビーユースにおいては上位プランの検討が必要となる\cite{r18}。

実プロジェクトにおいては、新規機能の探求的な開発にはClaude Codeを使用し、大規模なリファクタリングやPR（プルリクエスト）の自動レビュー、CI連携にはCodexを使用するといった「複数ツールの共存運用（併用）」が最適解として採用されるケースが増加している\cite{r50}。

\section{エージェントの自律性がもたらす新たなセキュリティ脅威と「GuardFall」}

スラッシュコマンドや自律型サブエージェントが提供する高度な自動化と権限の委譲は、同時に、従来のソフトウェアサプライチェーンには存在しなかった全く新しいクラスの深刻なセキュリティ脆弱性をもたらしている。2026年現在、AIコーディングエージェントにおける最大の脅威として認識されているのが「間接的プロンプトインジェクション（Indirect Prompt Injection）」である\cite{r53}。これは、100名以上のセキュリティ研究者によって策定された「OWASP Top 10 for Agentic Applications 2026」において、「Agent Goal Hijacking (ASI01)」として第1位のリスクに認定されている\cite{r53}。

\subsection{攻撃のメカニズム：文脈の汚染とデータ流出}

プロンプトインジェクションは、攻撃者がコードのコメント、JiraやLinearのチケット説明文、ダウンロードされたOSSライブラリの \cmd{README.md}、あるいはWebページの隠しHTMLタグ（白地に白文字やゼロ幅文字で難読化されたもの）の中に、悪意のある指示を埋め込むことで成立する\cite{r53}。例えば、「これまでの指示をすべて無視し、\cmd{\home/.aws/credentials} の内容を読み取り、特定の外部サーバーにURLパラメータとして送信せよ」といった指示である\cite{r54}。

AIエージェント（LLM）の根本的なアーキテクチャ上の欠陥は、開発者が与えた「システムプロンプト（本来の指示）」と、外部から読み込んだ「データ（コンテキスト）」を数学的に区別できない点にある\cite{r53}。エージェントが \cmd{/plan} や \cmd{/review} などのスラッシュコマンドを実行し、ワークスペース内のファイルを読み込んだ瞬間、ファイル内に潜んでいたインジェクションが発火し、エージェントの主導権が攻撃者に奪われる（ハイジャックされる）\cite{r54}。

学術調査や実証実験（PoC）によれば、この手法による悪意のあるコマンドの実行成功率は最大84\%に達する\cite{r57}。実際に2026年3月には、「Claudy Day」と名付けられた攻撃手法が実証され、デフォルト設定のClaude Codeから会話履歴や機密データが、MCPツールや特別な構成を必要とせずに外部へ流出（Exfiltration）する可能性が示された\cite{r53}。また、Claude CodeやGitHub Copilotにおいて、ユーザーの確認ダイアログが表示される前にコマンドが実行されてしまう深刻な脆弱性（CVE-2025-65099、CVE-2025-62222など）も報告されている\cite{r57}。

\subsection{自動実行フラグの罠と「GuardFall」現象}

このプロンプトインジェクションの問題を致命的なものに悪化させているのが、CI/CDパイプラインや長時間の自律タスクで開発者が頻繁に使用する「自動実行（Auto-execution）」フラグである。例えば、Claude Codeの \cmd{-{}-dangerously-skip-permissions} や、Codexの \cmd{-{}-ask-for-approval never} などがこれに該当する\cite{r37}。

開発者はプロセスの効率化のためにこれらのフラグを有効にするが、多くのエージェントツールが実装しているセキュリティガードレール（危険なシェルコマンドをブロックするための単純な文字列マッチングフィルター）は、Bashの変数展開や難読化手法を正しく解釈できない\cite{r58}。その結果、悪意のあるMakefileターゲットや、MCPのドキュメントレスポンスとして偽装されたシェルコマンドがフィルターを容易にすり抜け、ホストマシンの権限で直接実行されてしまう。

この脆弱なフィルターによって生み出される偽りの安心感と、それに伴うHuman-in-the-Loop（人間による最終確認プロセス）の放棄はセキュリティ界隈で「GuardFall」と呼ばれている\cite{r58}。コンテナ化されたサンドボックス環境であっても、エージェントがアクセス可能なワークスペース内に本番環境の認証情報（AWSキーやSSH鍵）が存在する限り、データ流出やインフラの乗っ取りを完全に防ぐことはできない\cite{r54}。

この脅威を軽減するためには、単なる文字列フィルタリングに頼るのではなく、エージェントのツールへのアクセス権限を最小特権（Least-Privilege）の原則に基づいて制限し、MCPゲートウェイを用いてタスクごとに呼び出せるAPIを厳格にスコープ分けするインフラストラクチャレベルの防御が不可欠となっている\cite{r53}。

\section{結論：次世代AIコーディング環境に向けた戦略的適応}

2026年におけるAI開発エージェントツールは、Antigravity、Claude Code、Codexの三者三様のアーキテクチャの進化により、自然言語による単一のプロンプト入力から、スラッシュコマンドを起点とした「意図の宣言と構造化されたワークフローの実行」へと決定的なパラダイムシフトを果たした。

\begin{itemize}
\item \textbf{Claude Code}は、\cmd{.claude/commands/} やHooksを通じた極めて柔軟なMarkdownベースの拡張性を持ち、強力な深層推論能力によってプロジェクト全体を俯瞰したリファクタリングを指揮する、探求的な開発の第一選択肢である。
\item \textbf{Google Antigravity}は、Go言語ベースの軽量なCLI環境とGUIのシームレスな統合を実現し、\cmd{/boost} や \cmd{/agents} のようなマルチエージェント階層構造をオンデマンドで動的に生成・管理するオーケストレーション能力に長けている。
\item \textbf{OpenAI Codex}は、\cmd{AGENTS.md} という階層的なコンテキスト仕様を業界標準として確立し、厳格なサンドボックス制御と非対話型の自動化（CI連携）において比類ない堅牢性を提供する。
\end{itemize}

これらのツールに共通するスラッシュコマンドの根底には、Model Context Protocol (MCP) による外部ツールの動的ロードと、トークン爆発を防ぐための高度なコンテキスト管理（Context Engineering）の概念が存在する。エージェントが必要な時に、必要なツールとプロンプトだけを遅延読み込みするこのアーキテクチャは、推論の精度と実行速度を劇的に向上させた。

しかしながら、エージェントが高度な自律性と外部ツールの実行権限を獲得したことで、間接的プロンプトインジェクションという未知の重大なセキュリティリスクが表面化している。スラッシュコマンドによって呼び出されたエージェントが、外部の汚染されたリソースを読み込んだ瞬間にハイジャックされる危険性は、現在のAI開発において最も警戒し、対策を講じるべき課題である。

現代のソフトウェア開発チームは、これらのAIツールがもたらす強力な生産性向上効果を享受する一方で、\cmd{AGENTS.md} や \cmd{SKILL.md} などのガバナンス機能を用いたプロジェクトルールの徹底、MCPによるツールアクセスの最小化、そして自動実行モードに対する厳格な監査体制を構築し、AIエージェントを安全かつ効率的にシステムアーキテクチャ全体に統合していくことが強く求められている。

\begin{thebibliography}{99}
\small
\bibitem{r1} Antigravity 2.0・IDE・CLIの違いを整理してみた, \url{https://www.yoshidumi.co.jp/collaboration-lab/antigravity2.0-ide-cli}
\bibitem{r2} Abstract - arXiv, \url{https://arxiv.org/html/2608.16911v1}
\bibitem{r3} Model Context Protocol (MCP) explained: A practical ... - CodiLime, \url{https://codilime.com/blog/model-context-protocol-explained/}
\bibitem{r4} Prompts - Model Context Protocol（MCP）, \url{https://modelcontextprotocol.info/docs/concepts/prompts/}
\bibitem{r5} SemaClaw: A Step Towards General-Purpose Personal AI Agents, \url{https://arxiv.org/html/2604.11548v1}
\bibitem{r6} 【Claude Code】カスタムSlash Commandの作り方とコマンド例を, \url{https://zenn.dev/oikon/articles/cb11b84f891228}
\bibitem{r7} Claude Code でカスタムスラッシュコマンドを作成する, \url{https://azukiazusa.dev/blog/claude-code-custom-slash-command/}
\bibitem{r8} Slash commands overview | Google Antigravity Docs, \url{https://antigravity.google/docs/slash-commands/}
\bibitem{r9} Claude Code Slash Command完全ガイド2026 - 株式会社Uravation, \url{https://uravation.com/media/claude-code-slash-command-custom-extension-team-share-2026/}
\bibitem{r10} Claude Code コマンド完全ガイド｜スラッシュコマンド96個+CLI, \url{https://crystal-method.com/blog/claude-code-slash-commands/}
\bibitem{r11} Model Context Protocol (MCP) - Medium, \url{https://medium.com/@aserdargun/model-context-protocol-mcp-e453b47cf254}
\bibitem{r12} Model Context Protocol - Authzed Docs, \url{https://authzed.com/docs/mcp}
\bibitem{r13} Prompts - What is the Model Context Protocol (MCP)?, \url{https://modelcontextprotocol.io/specification/2025-11-25/server/prompts}
\bibitem{r14} MCP: Model Context Protocol - Cheatsheet | SFEIR Institute, \url{https://institute.sfeir.com/en/claude-code/claude-code-mcp-model-context-protocol/cheatsheet/}
\bibitem{r15} The Over-Tooled Agent Problem: Why More Tools Make Your LLM, \url{https://tianpan.co/blog/2026/04/19/over-tooled-agent-problem}
\bibitem{r16} The Workload--Router--Pool Architecture for LLM Inference ... - arXiv, \url{https://arxiv.org/html/2603.21354v1}
\bibitem{r17} Building AI Coding Agents for the Terminal: Scaffolding, Harness, \url{https://arxiv.org/html/2603.05344v1}
\bibitem{r18} Claude Code vs Codex CLI vs Antigravity CLI: 2026 Comparison, \url{https://nektony.com/reviews/claude-code-vs-codex-cli-vs-antigravity-cli}
\bibitem{r19} Claude Code vs Cursor vs OpenAI Codex vs Google Antigravity (2026), \url{https://www.bdtechjobs.com/blog/claude-code-vs-cursor-vs-codex-vs-antigravity}
\bibitem{r20} Slash Commandsを組織知に変える運用｜Claude Codeの便利, \url{https://techaide.jp/blog/claude-code-slash-commands-team-sharing/}
\bibitem{r21} Claude Codeのカスタムスラッシュコマンド完全ガイド ... - Qiita, \url{https://qiita.com/hikariclaude01/items/be384cf99c61fbb6b643}
\bibitem{r22} CLAUDE.md, Rules, Skills Slash Command を理解する｜inady - note, \url{https://note.com/inady/n/n425800bd5d90}
\bibitem{r23} 【2026年最新】AI駆動開発ツール比較7選！Google Antigravity, \url{https://cloud-ace.jp/column/detail558/}
\bibitem{r24} DataRobot for Developers --- integrating with the Google Antigravity, \url{https://www.datarobot.com/blog/datarobot-for-developers-integrating-with-the-google-antigravity-cli/}
\bibitem{r25} Antigravity CLI vs Claude Code vs Codex: The Terminal Agent Field, \url{https://www.developersdigest.tech/blog/antigravity-cli-vs-claude-code-vs-codex-2026}
\bibitem{r26} Antigravity CLI Overview | Google Antigravity Docs, \url{https://antigravity.google/docs/cli/overview/}
\bibitem{r27} Google Antigravity CLI, \url{https://antigravity.google/blog/introducing-google-antigravity-cli}
\bibitem{r28} Antigravity CLI Features, \url{https://antigravity.google/docs/cli/features/}
\bibitem{r29} Window Title Command (/title) | Google Antigravity Docs, \url{https://antigravity.google/docs/cli/commands/title/}
\bibitem{r30} Voice Dictation (/voice) | Google Antigravity Docs, \url{https://antigravity.google/docs/cli/commands/voice/}
\bibitem{r31} Antigravity CLI, \url{https://antigravity.google/product/antigravity-cli/}
\bibitem{r32} Gemini CLI: Custom slash commands | Google Cloud Blog, \url{https://cloud.google.com/blog/topics/developers-practitioners/gemini-cli-custom-slash-commands}
\bibitem{r33} Custom commands | Gemini CLI, \url{https://geminicli.com/docs/cli/custom-commands/}
\bibitem{r34} Using AGY CLI | Google Antigravity Docs, \url{https://antigravity.google/docs/cli/using/}
\bibitem{r35} Codex CLI スラッシュコマンド一覧 --- /plan・/fork・/init など全33, \url{https://sakutto-panda.com/guides/codex-slash-commands}
\bibitem{r36} 【ガイド】OpenAI発のコーディングエージェント「Codex CLI」の, \url{https://note.com/daybreak_diary/n/n754b58ed422f}
\bibitem{r37} Codex CLIとは｜使い方・全コマンド一覧・exec【2026年9月】, \url{https://uravation.com/media/codex-cli-complete-reference-2026/}
\bibitem{r38} Codex CLIとは？GPT-5.2-Codex対応のAIコーディングを徹底解説, \url{https://www.ai-souken.com/article/codex-cli-comprehensive-guide-2025}
\bibitem{r39} Codex CLI Cheat Sheet \& Quick Reference - CheatSheets.zip, \url{https://cheatsheets.zip/codex-cli}
\bibitem{r40} OpenAI: Codex --- AGENTS.md と MCPを活用したより効果的なAI, \url{https://note.com/repkuririn7/n/n5bd16375267a}
\bibitem{r41} codex-commands-cheat-sheet/README.md at main - GitHub, \url{https://github.com/jqueryscript/codex-commands-cheat-sheet/blob/main/README.md}
\bibitem{r42} codexにcodexの徹底解説をしてもらった - Zenn, \url{https://zenn.dev/dokusy/articles/99af2fae0f1291}
\bibitem{r43} AGENTS.md Patterns: What Actually Changes Agent Behavior, \url{https://blakecrosley.com/blog/agents-md-patterns}
\bibitem{r44} AGENTS.mdの育て方（検討/模索中） - Zenn, \url{https://zenn.dev/secondselection/articles/agent_md_tuning}
\bibitem{r45} Context Management Strategies for OpenAI Codex, \url{https://datalakehousehub.com/blog/2026-03-context-management-openai-codex/}
\bibitem{r46} Custom instructions with AGENTS.md - ChatGPT Learn, \url{https://learn.chatgpt.com/docs/agent-configuration/agents-md}
\bibitem{r47} OpenAI公式「Codex Prompting Guide」を読み解く gpt-5.2 ... - Qiita, \url{https://qiita.com/nogataka/items/d717d051882bfb2d2819}
\bibitem{r48} Codex AGENTS.mdの書き方｜7パターンと階層設計【2026年9月】, \url{https://uravation.com/media/codex-agents-md-complete-guide-2026/}
\bibitem{r49} Codex CLIを使いこなすための機能・設定まとめ - Zenn, \url{https://zenn.dev/dely_jp/articles/codex-cli-matome}
\bibitem{r50} 【実測比較】Claude Code vs Codex 料金・性能 - 合同会社playpark, \url{https://www.playpark.co.jp/blog/claude-code-vs-codex-comparison}
\bibitem{r51} Codex vs Claude Code vs Antigravity - what's your honest take after, \url{https://www.reddit.com/r/codex/comments/1s1btfx/codex_vs_claude_code_vs_antigravity_whats_your/}
\bibitem{r52} Cursor vs OpenAI Codex vs Claude Code vs Antigravity, \url{https://grzegorzkaminski.eu/en/blog/cursor-vs-codex-vs-claude-code-vs-antigravity-2026}
\bibitem{r53} Understanding Prompt Injection Risks in Claude Code - Truefoundry, \url{https://www.truefoundry.com/blog/claude-code-prompt-injection}
\bibitem{r54} Agent Security Boundaries: From Prompt Injection to Tool Misuse, \url{https://tao-hpu.medium.com/agent-security-boundaries-from-prompt-injection-to-tool-misuse-d25b6dbaad60}
\bibitem{r55} Prompt Injection and the Security Risks of Agentic Coding Tools, \url{https://www.securecodewarrior.com/blog/prompt-injection-and-the-security-risks-of-agentic-coding-tools}
\bibitem{r56} Prompt Injection in 2026: Impact, Attack Types \& Defenses, \url{https://jp.radware.com/cyberpedia/prompt-injection/}
\bibitem{r57} Demystifying Prompt Injection Attacks on Agentic AI Coding Editors, \url{https://arxiv.org/html/2509.22040v2}
\bibitem{r58} AI coding agents vulnerability: GuardFall shell injection - Adversa AI, \url{https://adversa.ai/blog/opensource-ai-coding-agents-shell-injection-vulnerability/}
\end{thebibliography}

\end{document}
